\documentclass[10pt,a4paper]{amsart}
\usepackage[margin=19mm]{geometry}
\usepackage{amsmath,amssymb,mathtools}
\usepackage[scr=rsfs,cal=euler]{mathalfa}
\usepackage{xfrac}
\usepackage[unicode,hidelinks,bookmarks=false]{hyperref}
\newcommand{\ie}{i.e.\ }
\newcommand{\cf}{cf.\ }
\newcommand{\resp}{respectively\ }
\newcommand{\Subsection}{\subsection}
\providecommand{\nobreakdash}{\nobreak\hbox{-}}
\setlength{\emergencystretch}{2em}
\multlinegap0pt
\usepackage{bm}
\usepackage{bbm}
\usepackage{dsfont}
\usepackage{xspace}
\usepackage{cancel}
\usepackage{mathtools}
\usepackage{amsthm}
\makeatletter
\@ifundefined{Lemma}{\newtheorem{Lemma}{Lemma}}{}
\makeatother
\title[Global existence for the relativistic Vlasov-Poisson system ]{Global existence for the relativistic Vlasov-Poisson system in a two-dimensional bounded domain}
\author{Yanmin Mu, Dehua Wang}
\date{Source: arXiv:2511.06595v1 (CC BY 4.0)}
\begin{document}
\maketitle

\noindent\emph{Source and licence.} Global existence for the relativistic Vlasov-Poisson system in a two-dimensional bounded domain. Version arXiv:2511.06595v1 (CC BY 4.0), distributed under the Creative Commons Attribution 4.0 International licence, as recorded in the arXiv licence metadata for this exact version and shown on the abstract page https://arxiv.org/abs/2511.06595v1. Full rights notice: licenses/arxiv-2511.06595-CC-BY-4.0.txt.\par\smallskip

\noindent\textit{Editorial scope.} Counted unit: the Lemma whose source label is \texttt{L3} in the article (source0.tex lines 1411--1422, source label \texttt{L3}), delivered together with the article's own complete original proof of it (source0.tex lines 1423--1646, one \texttt{proof} environment of 224 lines). No mathematics is written, repaired or re-derived: statement and proof are the article's own text line for line. Required context retained uncounted: none. Every definition, display and label of those lines is retained unchanged. Omitted from this package and not redistributed: 1--1410, 1647--1899. Nothing inside a retained excerpt is omitted or altered: every retained body is delivered exactly as the article's lines are written, so no omission receipt of any kind accompanies this package. Citations: the retained excerpts carry no citation command at all, so no bibliography entry of the article is transcribed and no reference is redistributed. Reference adaptation: none was needed and none was made; every reference inside the delivered unit resolves inside it, so no unresolved reference or citation is produced. Printed slips kept literal: no printed word, symbol, hypothesis or number of the counted statement or of its proof was altered. Layout and class: the article's own class is \texttt{amsart} with packages including amsmath, amsfonts, amsthm, amssymb, graphicx, color, cite, enumitem, geometry, mathtools, hyperref, mathrsfs, titletoc, enumerate; the wrapper is the lane's pinned \texttt{amsart} editorial wrapper at 10pt on A4, which restates the theorem-like environments the retained text needs and adds the article's own macro definitions that the retained text needs (none), with extra wrapper packages (bm, bbm, dsfont, xspace, cancel, mathtools). The delivered numbering is the unit's own. Assets: no retained span contains a figure environment, a graphics-inclusion command, a PDF-page inclusion, a table or a TikZ picture; no image file is copied into this package, the package declares no asset dependency and no image file is redistributed with it. Licence: version arXiv:2511.06595v1 of this article is distributed under the Creative Commons Attribution 4.0 International licence, as recorded in the arXiv licence metadata for this exact version and shown on the abstract page https://arxiv.org/abs/2511.06595v1. A full rights notice is delivered at licenses/arxiv-2511.06595-CC-BY-4.0.txt. Characters: every retained excerpt is pure ASCII, so the delivered engine, XeLaTeX with its default Latin Modern text font, reports no missing character.
\begin{Lemma}\label{L3}
Assume further that  $\varphi(t,x)$ satisfies the following boundary value problem
\begin{align}
\Delta \varphi(t,x)&=\rho(t,x), \qquad (t,x)\in[0,T]\times\Omega,\nonumber\\
\varphi(t,x)&=0,\qquad \qquad x\in\partial\Omega.\nonumber
\end{align}
Then we have
\begin{align}
\Big|\frac{\partial \varphi}{\partial t}(t,\tilde{x})\Big|\leq C\tilde{x}_{1}\big(1+|\log\tilde{x}_{1}|+|\log\tilde{x}_{1}|^{2}\big),\nonumber
\end{align}
where $C>0$ depends only on $L=diam \Omega$ and $\|j\|_{\infty}$.
\end{Lemma}

\begin{proof}
Let's set $R=|\tilde{x}|\ll 1$, and do the following change of variables,
$$X=\frac{x}{R},\quad Y=\frac{y}{R},\quad \tilde{X}=\frac{\tilde{x}}{R}=(1,0).$$
The Green function for the domain $\Omega$ is denoted as $G$, therefore $\varphi$  can be expressed as
$$\varphi(t,x)=\int_{\Omega}\rho(t,y)G(x,y){\rm d}y,$$
then, letting $\Omega_{R}=\frac{\Omega}{R}$ and $H(X,Y):=G(RX,RY)$,   we obtain
\begin{align}
|\varphi_{t}(t,x)|&\leq\|j\|_{\infty}\int_{\Omega}|\nabla_{y}G(x,y)|{\rm d}y \nonumber\\
&=R\|j\|_{\infty}\int_{\Omega_{R}}|\nabla_{Y}H(X,Y)|{\rm d}Y.\nonumber
\end{align}
We therefore only need to prove that
\begin{align}
\int_{\Omega_{R}}|\nabla_{Y}H(X,Y)|{\rm d}Y\leq C\big(1+|\log R|+|\log R|^{2}\big).\label{v7.4}
\end{align}
The integration region must therefore be split into two parts, which are $|Y|\leq 4$ and $|Y|\geq 4$.

{\bf Case 1.}  When $|Y|\leq 4$, let $X=(X_{1},X_{2}),\,X^{*}=(-X_{1},X_{2}) $ is the reflextion of X with respect to the coordinate axis $\{X_{1}=0\}$.
First of all we define
$$\bar{H}(X,Y)=-\frac{1}{2\pi}\big(\log|Y-X|-\log|Y-X^{*}|\big),$$
which is the Green function for the half-space restricted to $\Omega_{R}\times\Omega_{R}$.

From a straightforward computation, it can be deduced that
\begin{align}
\int_{|Y|\leq 4}|\nabla_{Y}\bar{H}(\tilde{X},Y)|{\rm d}Y&\leq \frac{1}{2\pi}\int_{|Y|\leq 4}\big(\frac{1}{|Y-\tilde{X}|}+\frac{1}{|Y-\tilde{X}^{*}|}\big){\rm d}Y\nonumber\\
&\leq C. \nonumber
\end{align}

Let $W(X,Y)=H(X,Y)-\bar{H}(X,Y)$ and $\psi(X,Y)=\nabla_{Y}W(X,Y)$, then the function $\psi(X,Y)$ satisfies the following system:
\begin{align}
&\Delta_{X}\psi(X,Y)=0, \qquad X,Y\in\Omega_{R},\nonumber\\
&\psi(X,Y)=-\frac{1}{2\pi}\nabla_{Y}\big(\log|Y-X|-\log|Y-X^{*}|\big),\quad X\in\partial\Omega_{R}.\nonumber
\end{align}
%
If $\text{dist}(Y,\partial\Omega_{R})\geq 1$, then $|Y-X|\geq 1$ for all $X\in\partial\Omega_{R}$,  and we get
\begin{align}
|\psi(X,Y)|&=\frac{1}{2\pi}\big|\frac{Y-X}{|Y-X|^{2}}-\frac{Y-X^{*}}{|Y-X^{*}|^{2}}\big|\nonumber\\
&\leq C \big(\frac{1}{|Y-X|}+\frac{1}{|Y-X^{*}|}\big)\leq C\frac{1}{|Y-X|}\leq C.\nonumber
\end{align}
 According to the maximum principle, there is a constant $C>0$ uniformly with respect to $Y$, such that
 \begin{align}
|\psi(X,Y)|\leq C, \qquad  \forall X\in \Omega_{R}.\nonumber
\end{align}

Although the corresponding estimates  of $\psi(X,Y)$  become more intricate when the distance between $Y$ and $\partial\Omega_{R}$ is less than 1, the question can be resolved through the construction of a supersolution using the Poisson integral formula.
 Locating a boundary point $Y_{0}\in\partial\Omega_{R}$ such that $\text{dist}(Y,Y_{0})=\text{dist}(Y,\partial\Omega_{R})$ is a simple task. Then the following facts are true for $X\in \partial\Omega_{R}$.

If $|X-Y_{0}|\geq 2$, the triangle inequality  yields $|X-Y|\geq 1$ and then we get
\begin{align}
|\psi(X,Y)|\leq C\frac{1}{|X-Y|}\leq C. \nonumber
\end{align}

If $|X-Y_{0}|\leq R$, then we have
$$|\psi(X,Y)|\leq C\frac{1}{|X-Y|}.$$

If $R\leq|X-Y_{0}|\leq 2$, with the help of the Taylor theorem and the convexity of $\Omega$, we can deduce
$$|\psi(X,Y)|\leq \frac{C}{|X-Y_{0}|^{2}}.$$

To sum up,  we define
\begin{equation}\nonumber
\tilde{\psi}(X,Y)=\left\{\begin{aligned}
 &C, \qquad  if\,\, |X-Y_{0}|\geq 2, \\
&\frac{C}{|X-Y|}, \quad  if\,\, |X-Y_{0}|\geq R,\\
&\frac{C}{|X-Y_{0}|^{2}}, \quad  if\,\, R\leq|X-Y_{0}|\leq 2.\\
\end{aligned}  \right.
\end{equation}
Thus, for $|Y|\leq 4$ and $\text{dist}(Y,\partial\Omega_{R})\leq 1$, by means of the maximum principle, we have the following estimate, 
\begin{align}
|\psi(\tilde{X},Y)|\leq & C\int_{\partial\Omega_{R}}\frac{1}{|\tilde{X}-\xi|}|\tilde{\psi}(\xi,Y)|{\rm d}l
\leq C\int_{\partial\Omega_{R}}\frac{1}{1+|\xi|}|\tilde{\psi}(\xi,Y)|{\rm d}l\nonumber\\
=& C\int_{|\xi-Y_{0}|\leq R}\frac{1}{1+|\xi|}|\tilde{\psi}(\xi,Y)|{\rm d}l
+C\int_{R\leq|\xi-Y_{0}|\leq 2}\frac{1}{1+|\xi|}|\tilde{\psi}(\xi,Y)|{\rm d}l\nonumber\\
&+C\int_{|\xi-Y_{0}|\geq 2}\frac{1}{1+|\xi|}|\tilde{\psi}(\xi,Y)|{\rm d}l.\nonumber
\end{align}
Next, we estimate each item on the right end of the above equality,
\begin{align}
\int_{|\xi-Y_{0}|\geq 2}\frac{1}{1+|\xi|}|\tilde{\psi}(\xi,Y)|{\rm d}l&\leq C \int_{|\xi-Y_{0}|\geq 2}\frac{1}{1+|\xi|}{\rm d}l
\leq C |\log R|,\nonumber\\
\int_{R\leq|\xi-Y_{0}|\leq 2}\frac{1}{1+|\xi|}|\tilde{\psi}(\xi,Y)|{\rm d}l&\leq C \int_{R\leq|\xi-Y_{0}|\leq 2}\frac{1}{1+|\xi|}\cdot\frac{1}{|\xi-Y_{0}|^{2}}{\rm d}l
\leq C ,\nonumber
\end{align}
and leting $\eta=Y-Y_{0}$ and $|X-Y|=|(X-Y_{0})-\eta|$,   we have
\begin{align}
\int_{|\xi-Y_{0}|\leq R}\frac{1}{1+|\xi|}|\tilde{\psi}(\xi,Y)|{\rm d}l&\leq C \int_{|\xi-Y_{0}|\leq R}\frac{1}{1+|\xi|}\frac{1}{|\xi-Y|}{\rm d}l\nonumber\\
&=C\int_{|\xi-Y_{0}|\leq R}\frac{1}{1+|\xi|}\frac{1}{|\xi-Y_{0}-\eta|}{\rm d}l\nonumber\\
&\leq  C \int_{|\xi-Y_{0}|\leq R}\frac{1}{|\xi-Y_{0}|+|\eta|}{\rm d}l\nonumber\\
&\leq  C \int_{|\xi-Y_{0}|\leq R}\frac{1}{|\xi|+|\eta|}{\rm d}\xi
\leq C\log (1+\frac{R}{|\eta|})\nonumber\\
&\leq C+ C |\log R|+C|\log dist(Y,\partial\Omega_{R})|.\nonumber
\end{align}
Thus, together we obtain
\begin{align}
|\psi(\tilde{X},Y)|\leq C+ C |\log R|+C|\log dist(Y,\partial\Omega_{R})|,\nonumber
\end{align}
and hence,
\begin{align}
\int_{dist(Y,\partial\Omega_{R})\leq 1, |Y|\leq 4}|\psi(\tilde{X},Y)|{\rm d}Y&\leq C + C |\log R|\nonumber\\
&+C\int_{dist(Y,\partial\Omega_{R})\leq 1, |Y|\leq 4}|\log dist(Y,\partial\Omega_{R})|{\rm d}Y\nonumber\\
&\leq C+ C |\log R|.\nonumber
\end{align}
Therefore,  we deduce that
\begin{align}
\int_{ |Y|\leq 4}|\nabla_{Y}H(\tilde{X},Y)|{\rm d}Y&\leq \int_{ |Y|\leq 4}|\nabla_{Y}\bar{H}(\tilde{X},Y)|{\rm d}Y +\int_{ |Y|\leq 4}|\psi(\tilde{X},Y)|{\rm d}Y\nonumber\\
&\leq C+C|\log R|. \nonumber
\end{align}

{\bf Case 2.}  If $|Y|\geq 4$, we fix $Y=Y_{*}$, and rescale the  variables by $\eta=\frac{Y}{|Y_{*}|}, \zeta=\frac{X}{|Y_{*}|}$ such that 
$$\zeta,\eta\in \hat{\Omega}:=\frac{\bar{\Omega}_{R}}{|Y_{*}|}.$$
Define $g(\zeta,\eta)=H(|Y_{*}|\zeta,|Y_{*}|\eta)=H(X,Y)$. With the help of a change of variables, $g(\zeta,\eta)$ satisfies
$$\Delta_{\zeta}g(\zeta,\eta)=\frac{1}{|Y_{*}|}\delta(X-Y),$$
from the fact $\Delta_{X}H(X,Y)=\delta(X-Y)$.

Furthermore, let $\phi(\zeta,\eta)=|Y_{*}|\nabla_{\eta}g(\zeta,\eta)$, we obtain
\begin{align}
\Delta_{\zeta}\phi(\zeta,\eta)&=\nabla_{\eta}\delta(\zeta-\eta),\qquad   \zeta,\eta\in\hat{\Omega},\nonumber\\
\phi(\zeta,\eta)&=0,\qquad\qquad\qquad \zeta\in\partial\hat{\Omega}.\nonumber
\end{align}

Subsequently, we must further consider two cases depending on whether point $\eta$ is close to the boundary $\partial\hat{\Omega}$ or not.
\\
{\bf Case 2-1.} If $\text{dist}(\eta,\partial\hat{\Omega})\geq\frac{1}{10}$, $\psi(\zeta,\eta)$ is defined by
$$\psi(\zeta,\eta)=\phi(\zeta,\eta)+\frac{1}{2\pi}\nabla_{\eta}\log|\zeta-\eta|.$$
Then, the function $\psi(\zeta,\eta)$ satisfies the following system:
\begin{align}
\Delta_{\zeta}\psi(\zeta,\eta)&=0, \qquad \quad \zeta,\eta\in\hat{\Omega},\nonumber\\
\psi(\zeta,\eta)&=\frac{1}{2\pi}\nabla_{\eta}\log|\zeta-\eta|,\quad \zeta\in\partial\hat{\Omega}.\nonumber
\end{align}
The assumption yields $|\psi(\zeta,\eta)|\leq C$ for all $\zeta\in\partial\hat{\Omega}$. In the light of the maximum principle, we have
\begin{align}
|\psi(\zeta,\eta)|\leq C, \qquad \forall \zeta\in\partial\hat{\Omega}.\nonumber
\end{align}
%Now we consider the restricted region as $\{(\zeta,\eta)|\zeta\in\hat{\Omega}, |\zeta|\leq\frac{1}{2},|\eta|=1\}$,and use the boundary regularity theory for the Laplace operator, %then the following estimate holds
Given that we are now considering the confined region as $\{(\zeta,\eta)|\zeta\in\hat{\Omega}, |\zeta|\leq\frac{1}{2},|\eta|=1\}$ and applying the boundary regularity theory to the Laplace operator, the following estimate is valid,
\begin{align}
|\nabla_{\zeta}\psi(\zeta,\eta)|\leq C,  \qquad \forall \zeta\in\hat{\Omega}.\nonumber
\end{align}
Thus, for the restricted region it can be concluded that $|\nabla_{\zeta}\phi(\zeta,\eta)|\leq C,$  since $\phi(0,\eta)=0$. By the Taylor theorem, we get
\begin{align}
|\phi(\tilde{\zeta},\eta)|\leq C \text{dist}(\tilde{\zeta},\partial\hat{\Omega})=\frac{C}{|Y_{*}|}.\nonumber
\end{align}
Therefore, by the above range of variable substitution, we finally obtain that
$$|\nabla_{Y}H(\tilde{X},Y_{*})|\leq\frac{C}{|Y_{*}|^{3}}.$$
%
{\bf Case 2-2.} If $\text{dist}(\eta,\partial\hat{\Omega})\leq\frac{1}{10}$, we fix $\eta=\eta_{0}$ and define $\bar{\eta}_{0}$ as the boundary point closest to $\eta_{0}$.
 Let $\zeta^{*}$ as the reflection point
with respect to the tangent line at $\bar{\eta}_{0}$, namely,
\begin{align}
\zeta^{*}=\zeta+2\text{dist}(\zeta,\{\eta\in \mathbb{R}^{2},(\eta-\bar{\eta}_{0})\cdot n(\bar{\eta}_{0})=0\})n(\bar{\eta}_{0}),\nonumber
\end{align}
where $n(\bar{\eta}_{0})$ is the outward normal vector. The functions $\bar{g}(\zeta,\eta),\omega(\zeta,\eta)$ are defined as follows,
\begin{align}
\bar{g}(\zeta,\eta)&=|Y_{*}|\nabla_{\eta}g(\zeta,\eta),\nonumber\\
\omega(\zeta,\eta)&=\bar{g}(\zeta,\eta)+\frac{1}{2\pi}\nabla_{\eta}\big(\log|\zeta-\eta|-\log|\zeta^{*}-\eta|\big).\nonumber
\end{align}
Then we have
\begin{align}
&\Delta_{\zeta}\omega(\zeta,\eta)=0, \qquad \qquad\qquad \zeta,\eta\in\hat{\Omega},\nonumber\\
&\omega(\zeta,\eta)=\frac{1}{2\pi}\big(\frac{\zeta-\eta}{|\zeta-\eta|^{2}}-\frac{\zeta^{*}-\eta}{|\zeta^{*}-\eta|^{2}}\big),\quad \zeta\in\partial\hat{\Omega}.\nonumber
\end{align}
Now, for $\zeta\in\partial\hat{\Omega}$, by the triangle inequality, the following estimate holds

\begin{equation}\nonumber
|\omega(\zeta,\eta_{0})|\leq\left\{\begin{aligned}
 &C, \qquad  if\,\, |\zeta-\bar{\eta}_{0}|\geq \frac{1}{8}, \\
&\frac{C}{|\zeta-\bar{\eta}_{0}|+|\eta_{0}-\bar{\eta}_{0}|}, \quad  if\,\, |\zeta-\bar{\eta}_{0}|\leq \frac{1}{8}.\\
\end{aligned}  \right.
\end{equation}

We denote $D=\text{dist}(\eta_{0},\partial\hat{\Omega})=|\eta_{0}-\bar{\eta}_{0}|$, and apply the Poisson kernel formula to obtain, for $\zeta\in \mathbf{B}_{\frac{1}{|Y_{*}|}}(0)\cap\hat{\Omega}$,
\begin{align}
|\omega(\zeta,\eta_{0})|&\leq\int_{\partial\hat{\Omega}}\frac{1}{|\zeta-\beta|}|\omega(\beta,\eta_{0})|{\rm d}l\nonumber\\
&\leq\int_{|\beta-\bar{\eta}_{0}|\leq\frac{1}{8}}\frac{1}{|\zeta-\beta|}\cdot\frac{C}{|\zeta-\bar{\eta}_{0}|+|\eta_{0}-\bar{\eta}_{0}|}{\rm d}l+
\int_{|\beta-\bar{\eta}_{0}|\geq\frac{1}{8}}\frac{C}{|\zeta-\beta|}{\rm d}l\nonumber\\
&=:M_{1}+M_{2}\nonumber.
\end{align}

First, the estimation of $M_{1}$ is performed  as follows, since
$$|\eta_{0}|=1,\,\, |\eta_{0}-\bar{\eta}_{0}|\leq\frac{1}{10},\,\, |\zeta|\leq \frac{1}{4},$$
we have
$$|\zeta-\bar{\eta}_{0}|\geq\frac{13}{20},$$
then, for $\beta\in \partial\hat{\Omega}$ and $|\beta-\bar{\eta}_{0}|\leq\frac{1}{8}$, the triangle inequality yields
$$|\zeta-\beta|\geq\frac{21}{40}.$$
Thus,
\begin{align}
M_{1}\leq C\int_{0}^{\frac{1}{8}}\frac{1}{\mu+D}{\rm d}\mu\leq C|\log D|.\nonumber
\end{align}

Second, we estimate  $M_{2}$.  For $\beta\in \partial\hat{\Omega}$ and $|\beta-\bar{\eta}_{0}|\geq\frac{1}{8}$, due to
$$|\tilde{\zeta}|=|\frac{|\tilde{X}|}{|Y_{*}|}|=\frac{1}{|Y_{*}|},\,\,\text{dist}(\tilde{\zeta},\partial\hat{\Omega})=\frac{1}{|Y_{*}|},\,\,
|\tilde{\zeta}-|\leq\frac{1}{10},\,\, |\tilde{\zeta}-\beta|\geq \frac{1}{|Y_{*}|} ,$$
we get,
$$|\tilde{\zeta}-\beta|\geq \frac{1}{3}\big(\frac{1}{|Y_{*}|}+|\beta|\big),$$
therefore,
\begin{align}
M_{2}&\leq C\int_{|\beta-\bar{\eta}_{0}|\geq\frac{1}{8}}\frac{|Y_{*}|}{1+|Y_{*}||\beta|}{\rm d}l\nonumber\\
&\leq C|Y_{*}|\int^{\frac{C}{|Y_{*}|R}}_{0}\frac{1}{1+|Y_{*}|\mu}{\rm d}\mu\nonumber\\
&\leq C |\log R|.\nonumber
\end{align}

Finally, we obtain
\begin{align}
|\omega(\tilde{\zeta},\eta_{0})|&\leq C\big(|\log R|+|\log D|\big),\nonumber\\
|\bar{g}(\tilde{\zeta},\eta_{0})|&\leq C\big(1+|\log R|+|\log D|\big),\nonumber
\end{align}
thus, by the above range of variable substitution, we have
\begin{align}
|\nabla_{Y}H(\tilde{X},Y_{*})|\leq C\big(\frac{1}{|Y_{*}|^{2}}+\frac{1}{|Y_{*}|^{2}}|\log R|+\frac{1}{|Y_{*}|^{2}}|\log {\rm dist}(\frac{Y_{*}}{|Y_{*}|}.\partial\hat{\Omega}) |\big).\nonumber
\end{align}
Therefore,  if $|Y|\geq 4$,
\begin{align}
\int_{4\leq |Y|\leq \frac{L}{R}}|\nabla_{Y}H(\tilde{X},Y)|{\rm d}^{2}Y
 \leq C & \int_{4\leq |Y|\leq \frac{L}{R}}\frac{1}{|Y|^{2}}\big(1+|\log R| \nonumber\\
&+|\log dist(\frac{Y}{|Y|},\partial\hat{\Omega})|\big) {\rm d}^{2}Y \nonumber\\
\leq C & \big(|\log R|+|\log R|^{2}\big),\nonumber
\end{align}
where the following integral estimation has been used,
\begin{align}
\int_{4\leq |Y|\leq \frac{L}{R}}&\frac{1}{|Y|^{2}}|\log dist(\frac{Y}{|Y|},\partial\hat{\Omega})|{\rm d}^{2}Y \nonumber\\
&\leq\sum^{|\log_{2}R|}_{n=0}\int_{4\cdot 2^{n}\leq |Y|\leq 4\cdot 2^{n+1} }\frac{1}{|Y|^{2}}|\log dist(\frac{Y}{|Y|},\partial\hat{\Omega})|{\rm d}^{2}Y \nonumber\\
&\leq\sum^{|\log_{2}R|}_{n=0}\int_{4\leq |Z|\leq 8 }\frac{1}{|Z|^{2}}|\log dist(\frac{Z}{|Z|},\frac{\partial\hat{\Omega}}{2^{n}})|{\rm d}^{2}Z \nonumber\\
&\leq \frac{1}{16}\sum^{|\log_{2}R|}_{n=0}\int_{4\leq |Z|\leq 8 }|\log dist(\frac{Z}{|Z|},\frac{\partial\hat{\Omega}}{2^{n}})|{\rm d}^{2}Z \nonumber\\
&\leq C|\log R|.\nonumber
\end{align}
To this extent, we have completed the proof of this lemma.
\end{proof}

\end{document}
